\documentclass[11pt]{article}
\input{packages.tex}
\usepackage{soul}
\usepackage{multicol}
\DeclareMathOperator{\sign}{sign}
\DeclareMathOperator{\Imm}{Im}
\DeclareMathOperator{\rg}{rg}

% Настройка колонтитулов
\pagestyle{fancy}
\fancyhf{} % Очистить все поля
\fancyhead[L]{Сделаем вид, что это не копипаст презентации :)}
\fancyhead[R]{\href{https://t.me/lost_luna}{Автор}}
\fancyfoot[C]{\thepage}               % Справа внизу

\setlength{\headheight}{13.6pt}

\hypersetup{
    colorlinks=true,
    linkcolor=blue,
    citecolor=blue,
    filecolor=blue,
    urlcolor=blue,
}

\begin{document}

\begin{titlepage}
    \centering
    \vspace*{\fill}  % заполняет пространство сверху

    {\Huge Лекция 8. Грамматики. Машины Тьюринга.
    \par}

    \vspace*{\fill}  % заполняет пространство снизу
\end{titlepage}

\setcounter{page}{1}


\section{Полезный частный случай}

\begin{statement}
    Пусть $L \subseteq \{a\}^*$, то есть $L$ --- язык в однобуквенном алфавите. Тогда $L$ --- регулярный $\Longleftrightarrow$ он представим в виде объединения \underline{конечного}
    числа арифметических прогрессий, где под арифметической прогрессией понимается регулярный язык вида \[L_{nm} = \{a^n \cdot (a^m)^t, \ t \in \mathbb{N}_0\}\]
\end{statement}

\begin{proof}
    ($\Longleftarrow$): $L_{nm} = a^m \cdot (a^n)^*$ --- регулярное выражение, а в силу того, что таких арифметических прогрессий конечное число, мы можем пронумеровать их 
    и составить из них регулярное выражение: \[L = L_{n_1 m_1} + L_{n_2 m_2} + \dots + L_{n_k m_k} \text{ --- регулярное выражение}\] Следовательно, язык $L$ регулярный.\\
    \leavevmode
    ($\Longrightarrow$): пусть $L$ --- регулярный язык, тогда по нему можно построить ПДКА. Возьмем произвольное слово из языка, пусть ПДКА принимает это слово в некотором
    финальном состоянии $q_f$. Заметим, что $q_f$ либо лежит в каком-нибудь цикле, либо не лежит.\\
    Пусть в ПДКА $\left|Q\right| + 1$ состояний, тогда возможные простые циклы имеют длину не больше $\left|Q\right|$. Теперь рассмотрим кратчайшие пути к финальным состояниям
    из начального (то есть те пути, которые не посещают одну и ту же вершину два раза): таких путей также будет не более $\left|Q\right| + 1$. В множестве арифметических прогрессий
    $L_{nm}: \ m \leq \left|Q\right|$ содержатся все арифметические прогрессии, в объединении которых лежит язык $L$, так как каждое слово из языка лежит в некоторой 
    арифметической прогрессии, которая, в свою очередь, целиком лежит в языке $L$ (это следует из леммы о разрастании). Покажем, что подоходящих нам арифметических прогрессий,
    на самом деле, конечное число: зафиксируем $m \leq \left|Q\right|$, рассмотрим следующие арифметические прогрессии: \[a^{n_1} (a^m)^t, \ a^{n_2} (a^m)^t, \dots, a^{n_k} (a^m)^t\]
    Заметим, что как только $k > m$, среди прогрессий найдутся хотя бы $k-m$ вложенных (так как существует всего $m$ остатков от деления на $m$, и при $k > m$, по принципу Дирихле, хотя бы 
    $k - m$ арифметических прогрессий будут иметь неуникальные остатки), следовательно, есть не более $m$ арифметических прогрессий. Отсюда легко получаем, что число арифметических 
    прогрессий, в объединении которых лежит язык $L$, не превосходит $(\left|Q\right| + 1) \times (\left|Q\right| + 1)$, то есть, это число конечно.

\end{proof}

\section{Грамматики}

\begin{definition}
    Грамматика $G$ --- это набор из четырех множеств $(N, T, P, S)$, где 
    \begin{itemize}
        \item $N$ --- алфавит нетерминальных символов
        \item $T$ --- алфавит терминальных символов (символы из алфавита формального языка), $T \cap N = \varnothing$
        \item $P$ --- конечное множества правил вида $\alpha \to \beta$, где $\alpha \in (N \cup T)^* N (N \cup T)^*, \ \beta \in (N \cup T)^*$
        \item $S \in N$  --- начальный знак (аксиома) грамматики 
    \end{itemize}
\end{definition}

\begin{example}
    Зададим с помощью грамматики наш любимый язык $a^nb^n$: \[P = \{ S \to \varepsilon, \ S \to ab, \ S \to aSb\}, \quad N = \{S\}, \quad T = \{a, b\}\]
    Если последовательно применять правила, пока не получится цепочка, состоящая только из терминальных символов, то получатся только слова вида $a^n b^n$.
    \[S \to a S b \to aa S bb \to aaabbb = a^3 b^3\]
\end{example}

\begin{remark}
    Грамматика задает язык $L$, если любые слова из $L$ можно вывести из $S$, применяя правила из $P$ и только их.
\end{remark}

\begin{example}
    Построим грамматику по ДКА:
    \begin{figure}[H]
            \centering
            \includegraphics[width=0.5\textwidth]{pictures/Грамматика по ДКА.png}
            \caption{вот по такому ДКА}
        \end{figure}
    \leavevmode
    Поставим в соответствие каждому состоянию некоторый нетерминал, и построим множество правил: 
    \[P = \{S \to aA, \ A \to bA, \ A \to aB, \ B \to bS, \ B \to aC, \ B \to \varepsilon, \ C \to \varepsilon\}\]
    Например, слово $"aa"$ мы можем получить следующим образом: $S \to aA \to aaB \to aa$.
\end{example}

\begin{definition}
    Типы грамматик (иерархия Хомского):
    \begin{itemize}
        \item Грамматику, для которой нет дополнительных ограничений, называют грамматикой \textbf{без ограничений} (типа 0). (то есть в правилах нетерминалы и терминалы 
        могут стоять как угодно, а не только один нетерминал слева, как в рассматриваемых нами грамматиках)\\ 
        \textcolor{gray}{Практического применения не имеют, так как слишком сложная структура. Распознается МТ.}
        \item Если каждое правило грамматики кроме $S \to \varepsilon$ имеет вид $\alpha \to \beta$, где $|\alpha| \leq |\beta|$, и в случае, если правило $S \to \varepsilon \in P$, 
        то символ $S$ не встречается в правых частях правил, то грамматику называют контекстно-зависимой или неукорачивающей (типа 1)\\
        \textcolor{gray}{Используются для анализа естественных языков.}
        \item Если каждое правило грамматики имеет вид $A \to \beta$, где $A \in N$, $\beta \in (N \cup T)^*$, то её называют контекстно-свободной (типа 2)\\
        \textcolor{gray}{Широко применяются для описания синтаксиса компьютерных языков. Распознаются при помощи конечного автомата со стеком.}
        \item Если каждое правило грамматики имеет вид либо $A \to xB$ , либо $A \to x$, где $A, B \in N$, $x \in T^*$, то грамматику называют регулярной (типа 3)
        \textcolor{gray}{Совпадает с множеством регулярных. Используется для ассемблера, командных процессов и др.}
    \end{itemize}
\end{definition}


\section{Машины Тьюринга}

\begin{definition}
    Машина Тьюринга (МТ) — это набор $(A, Q, Q_f, \delta, q_0, \Lambda)$, где 
    \begin{itemize}
        \item $A$ --- конечное множество, называемое алфавитом МТ
        \item $Q$ --- конечное множество, называемое множеством состояний МТ: $A \cap Q = \varnothing$
        \item $Q_f$ --- множество финальных состояний МТ: $Q_f \subseteq Q$
        \item $\delta$ — функция (таблица) переходов: $\delta: A \times Q \to A \times \{-1, 0, 1\} \times Q$
        \item $q_0$ --- начальное состояние: $q_0 \in Q$
        \item $\Lambda$ --- пустой символ: $\Lambda \in A$
    \end{itemize}
\end{definition}

\begin{remark}
    Будем считать, что задача является алгоритмически разрешимой $\Longleftrightarrow$ по ней можно построить машину Тьюринга.
\end{remark}

\begin{remark}
    Будем считать, что машина Тьюринга завершает работу, как только она попадает в финальное состояние. При этом, если машина Тьюринга попала в финальное состояние,
    это не значит, что она приняла слово.
\end{remark} \leavevmode 
\\
Работа машины Тьюринга выглядит следующим образом: у нас имеется лента, она же память (бесконечная), разделенная на ячейки. В самом начале управляющее устройство машины 
Тьюринга смотрит на какую-то одну ячейку, при этом устройство находится в состоянии $q_0$. На вход МТ подается слово $x \in A^*$. Например, $A = \{a, b, \Lambda\}, \ x = aba$
(считаем, что в самом начале управляющее устройство указывает на первый символ поданного на вход слова). Мы работаем с детерминированными машинами Тьюринга, то есть, 
по состоянию и символу мы можем однозначно определить переход.
\begin{figure}[H]
    \centering
    \includegraphics[width=0.6\textwidth]{pictures/Начальное состояние МТ.png}
    \caption{начальное состояние МТ}
\end{figure} \leavevmode \\
Пусть $\delta(a, q_0) = b \times \{+1\} \times q_1$. Это значит, что символ $a$ поменяется на $b$, управляющее устройство сдвинется вправо и перейдет в состояние $q_1$.
\begin{figure}[H]
    \centering
    \includegraphics[width=0.6\textwidth]{pictures/После перехода.png}
    \caption{после перехода}
\end{figure}

\begin{statement}
    Множество всех машин Тьюринга счётно. (Любую машину Тьюринга можно представить конечным словом из 0 и 1, определив способ кодирования состояний и символов алфавита).
\end{statement}

\begin{example}
    Построим машину Тьюринга и закодируем ее: пусть \[A = \{a, b, c, \Lambda\}, \quad Q = \{q_0, q_1\}, \quad Q_f = \{q_1\}\]
    \[
    \delta: \quad \begin{array}{c|c|c|c|c}
        & a & b & c & \Lambda \\
        \hline
        q_0 & a, +1, q_1 & b, +1, q_1 & b, -1, q_1 & \Lambda, 0, q_1 \\
        \hline
        q_1 & & & &
    \end{array}
    \]
    Будем кодировать каждое слово четырьмя битами, чтобы обозначить это, запишем 4 единицы (до первого нуля). 0 --- разделяющий символ, его пропускаем. Теперь закодируем 
    алфавит: $\Lambda - 0000, \ a - 0001, \ b - 0010, \ c - 0011$. Для удобства символы алфавита будут начинаться с 0, а символы состояний --- с 1. Договоримся, что 0
    после множества состояний отделяет $Q$ от $Q_f$, также 0 будет отделять $Q_f$ от переходов. 
    \begin{figure}[H]
        \centering
        \includegraphics[width=0.6\textwidth]{pictures/Кодировка МТ.png}
        \caption{вот так закодируем}
    \end{figure} \leavevmode \\
    Перечислять переходы будем следующим образом: первые 4 бита --- символ алфавита, следующие 4 бита --- состояние, далее 4 бита --- новый символ алфавита, потом 2 бита (тут всегда 2 бита) на 
    $\pm 1, 0$, и 4 бита на новое состояние. (Почти везде присутствует магическое число 4, потому что мы изначально выбрали именно 4 бита для кодировки всех слов).
    \begin{figure}[H]
        \centering
        \includegraphics[width=0.6\textwidth]{pictures/Кодировка перехода.png}
        \caption{вот так закодируем переход}
    \end{figure} \leavevmode \\
    Ещё для удобства любое количество нулей в начале будем отбрасывать. \\
    Таким образом, любую машину Тьюринга мы можем закодировать конечным словом из 0 и 1, а значит, всего машин Тьюринга счётное количество.
\end{example}

\begin{remark}
    Так как языков несчётное количество, значит, не все языки алгоритмически разрешимы.
\end{remark}

\begin{remark}
    Другие формализации МТ (в этом семестре эти МТ эквивалентны, то есть задают одно и то же множество алгоритмически разрешимых задач):
    \begin{itemize}
        \item $k$-ленточная МТ 
        \item МТ с лентой, бесконечной в одну сторону
        \item МТ с лентой, бесконечной в обе стороны
        \item детерминированная / недетерминированная МТ
    \end{itemize}
\end{remark}
\end{document}